%% ma: L10-B14-0023
%% lop: 10
%% chuong: 5
%% bai: 14
%% dang: 10.14.6
%% loai: TN4
%% muc: VD
%% dap_an: "A"
%% nguon: "chuyen-de-cac-so-dac-trung-cua-mau-so-lieu-khong-ghep-nhom-toan-10-chuong-trinh-moi.pdf (GV. Phan Nhật Linh) (P04), Bài 03 Dạng 1 Phần I Câu 15"
%% kien_thuc: "Lớp 10 Bài 13 (tứ phân vị), Bài 14 (khoảng tứ phân vị)"
%% trang_thai: chinh-thuc
%% ghi_chu: "Đơn vị gốc ghi kw, đã sửa thành kWh. Tự duyệt theo yêu cầu giáo viên 22/10/2026."
\begin{ex}
Mẫu số liệu sau cho biết lượng điện tiêu thụ (đơn vị: kWh) hằng tháng của gia đình bạn An trong năm 2021:
\begin{center}
\begin{tabular}{|c|c|c|c|c|c|c|c|c|c|c|c|}\hline
$163$ & $165$ & $159$ & $172$ & $167$ & $168$ & $170$ & $161$ & $164$ & $174$ & $170$ & $166$\\\hline
\end{tabular}
\end{center}
Trong năm 2022, nhà bạn An giảm mức tiêu thụ điện mỗi tháng $10$ kWh (so với cùng tháng năm 2021). Gọi $\Delta_Q$, $\Delta'_Q$ lần lượt là khoảng tứ phân vị của mẫu số liệu tiêu thụ điện năm 2021 và năm 2022. Đẳng thức nào sau đây đúng?
\choice{\True $\Delta_Q=\Delta'_Q$}{$\Delta'_Q=\Delta_Q-10$}{$\Delta_Q=\Delta'_Q-10$}{$\Delta'_Q=\Delta_Q-20$}
\loigiai{Năm 2021: sắp xếp $159$; $161$; $163$; $164$; $165$; $166$; $167$; $168$; $170$; $170$; $172$; $174$, $Q_1=163{,}5$, $Q_3=170$, $\Delta_Q=6{,}5$.

Năm 2022 mọi giá trị giảm $10$ nên $Q_1$, $Q_3$ cùng giảm $10$: $\Delta'_Q=160-153{,}5=6{,}5=\Delta_Q$. Chọn A.}
\end{ex}
